\chapter{Behavioural Evidence: Where Rational Choice Breaks}\label{ch:b2:behavioural-evidence}

Eighteen laboratory experiments published in two leading economics journals between 2011 and 2014 were
% ledger: b2:camerer2016
run again, with new participants, a plan of analysis posted in advance and samples large enough to detect
the original effect nine times out of ten. Each dot in the chart below is one of them: the original
effect on the horizontal axis, the replication's on the vertical. Had every effect been reproduced
exactly, the dots would lie on the diagonal. Most lie below it.

\begin{openingchart}{Original against replication effect sizes for 18 laboratory experiments in
  economics, each measured as a correlation coefficient $r$. Filled dots: the replication found a
  significant effect in the original direction; open dots: it did not. The diagonal marks a replication
  equal to the original; the dotted line, an effect of zero.}{Camerer et al.\ (2016), Science,
  Table~S1 of the supplementary materials}
\begin{tikzpicture}
\begin{axis}[omrecord, grid=none, xmin=0, xmax=1, ymin=-0.2, ymax=1, xtick={0,0.2,0.4,0.6,0.8,1},
    ytick={-0.2,0,0.2,0.4,0.6,0.8,1}, xlabel={Original effect size ($r$)},
    ylabel={Replication effect size ($r$)}, width=0.62\linewidth, height=0.62\linewidth,
    legend style={at={(0.03,0.97)}, anchor=north west, font=\footnotesize}]
  \addplot[only marks, mark=*, mark size=2.2pt, omDef] table[x=orig_r, y=rep_r]
    {figdata/bachelor-2/05-behavioural-evidence/out/replication_yes.csv};
  \addplot[only marks, mark=o, mark size=2.4pt, omThm, thick] table[x=orig_r, y=rep_r]
    {figdata/bachelor-2/05-behavioural-evidence/out/replication_no.csv};
  \legend{Replicated, Not replicated}
  \draw[black!60] (axis cs:0,0) -- (axis cs:1,1);
  \draw[black!60, densely dotted, thick] (axis cs:0,0) -- (axis cs:1,0);
\end{axis}
\end{tikzpicture}
\end{openingchart}

% figdata: replication_summary (n 18, n_replicated 11, n_smaller 15, n_negative 3)
Eleven of the 18 replications found a significant effect in the original direction. Fifteen found a
smaller effect than the original, and three found an effect of the opposite sign. Most of the 18 studies
are not about the choices of this chapter, but the record sets its standard. Expected utility (\cref{ch:b2:expected-utility}) and exponential
discounting (\cref{ch:b2:intertemporal-choice}) describe a chooser whose preferences over final outcomes
are fixed and who uses every piece of information that is free. A large body of evidence finds choices
that these models do not predict. This chapter asks which of those findings hold up when they are
repeated, when stakes rise and when the people choosing have experience, and it states the best-known
alternative model, \omterm{def:b2:behavioural-evidence:prospect-theory}{prospect theory}, in the same frame as every other model of the book: what it assumes,
what it predicts, how it was tested and where it breaks.

\begin{recall}
From \cref{ch:b2:expected-utility}: the \omterm{def:b2:expected-utility:independence-axiom}{independence axiom}, the Allais pairs and the common-ratio
pattern (\cref{prop:b2:expected-utility:common-ratio}), and the calibration bound
(\cref{prop:b2:expected-utility:calibration}): a person who turns down a small 50--50 gamble at every
level of wealth must turn down very large ones. From \cref{ch:b2:intertemporal-choice}: \omterm{def:b2:intertemporal-choice:present-bias}{present bias}
(\cref{def:b2:intertemporal-choice:present-bias}), the quasi-hyperbolic factor $\varphi$, and the field
evidence on gym contracts and on effort. From the school volumes: a correlation coefficient, a standard
error and a significance test.
\end{recall}

\begin{notation}[Symbols of this chapter]\label{not:b2:behavioural-evidence:symbols}
Outcomes are $z$; $z_0$ is a \omterm{def:b2:behavioural-evidence:reference-point}{reference point}. $v$ is the \omterm{def:b2:behavioural-evidence:value-function}{value function} of \omterm{def:b2:behavioural-evidence:prospect-theory}{prospect theory}, $\kappa$ the
loss-aversion coefficient and $\eta$ the curvature of $v$; $\omega(p)$ is a \omterm{def:b2:behavioural-evidence:weighting}{probability weighting
function} with parameter $d$. In the model of defaults, $b_i$ is person $i$'s net benefit from joining a
plan and $\chi$ a switching cost; in the model of attention, $\tau$ is a tax and $\theta$ the weight
given to it. Effect sizes are correlation coefficients $r$. The literature writes the loss-aversion
coefficient as $\lambda$; this book keeps $\lambda$ for a Lagrange multiplier.
\end{notation}

\section{What a replication counts}\label{sec:b2:behavioural-evidence:replication}

A study reports an estimate, a standard error and, usually, a test of the hypothesis that the effect is
zero. A \emph{replication} runs the same design on new data. Two quantities summarise a set of
replications.

\begin{definition}[Replication rate]\label{def:b2:behavioural-evidence:replication-rate}
The \emph{replication rate}\index{replication rate} of a set of studies is the share of their
replications that find an effect significant at the chosen level (usually 5 per cent) in the same
direction as the original. The \emph{relative effect size} of a replication is its estimate divided by
the original estimate.
\end{definition}

The \omterm{def:b2:behavioural-evidence:replication-rate}{replication rate} depends on the power of the replications as well as on the truth of the effects: a
replication with power $0.9$ misses a true effect of the original size one time in ten. That is why the
18 replications of the opening chart were designed with power of at least 90 per cent. If every original
% ledger: b2:camerer2016
effect were true and accurately estimated, about 92 per cent of them should have replicated (the mean
power); 61.1 per cent did. The mean effect size fell from $r = 0.474$ in the originals to $0.279$ in the
replications, and the mean relative effect size was 65.9 per cent.

Why should effects shrink? The simplest reason needs no error by anyone. Journals and authors publish
significant results more readily than null ones. If a true effect is small and the samples modest, the
estimates that pass the significance filter are the ones that happened to come out large, so the
published estimate exaggerates the effect, while a replication, which is reported whatever it finds, is
unbiased.

\begin{proposition}[The significance filter exaggerates]\label{prop:b2:behavioural-evidence:filter}
Let an estimate be $\hat r = r + \sigma u$ with $r > 0$, $\sigma > 0$ and $u$ standard normal. Conditional
on passing the one-sided test $\hat r > 1.96\,\sigma$, the expected estimate exceeds $r$. A replication
$\hat r' = r + \sigma u'$, with $u'$ independent of $u$, has expectation $r$ whatever the first estimate
was.
\end{proposition}

\begin{proof}
$\mathbb{E}[\hat r \mid \hat r > c] = r + \sigma\,\mathbb{E}[u \mid u > (c - r)/\sigma]$, and the mean of
a standard normal variable above any threshold is positive. The replication's error $u'$ is independent
of the event that selected the original, so $\mathbb{E}[\hat r' \mid \hat r > c] = r$.
\end{proof}

The listing below simulates the proposition. A true effect of $r = 0.15$ is studied in 20,000 samples of
100 people; only the significant estimates are published; each published study is then replicated once
with the same sample size.

\omcode{code/bachelor-2/05-behavioural-evidence/listings/winners_curse.py}{1}{19}{Publication of
significant results only, and the replications that follow.}\label{lst:b2:behavioural-evidence:curse}

\omoutput{code/bachelor-2/05-behavioural-evidence/outputs/winners_curse.txt}{The published estimates
are inflated; the replications are not.}\label{out:b2:behavioural-evidence:curse}

% code: listings/winners_curse.py
Only about a third of the studies of the true effect pass the filter, and those that do report a mean
effect of 0.262, three quarters larger than the truth. Their replications average 0.149, the true value,
and only about a third of them are significant, so the field would record a \omterm{def:b2:behavioural-evidence:replication-rate}{replication rate} of about
one third for an effect that is real. The relative effect size is 0.568. A low \omterm{def:b2:behavioural-evidence:replication-rate}{replication rate} and
shrinking effects are therefore what a literature of small true effects and low power produces; they do
not show that the original effects were zero. What they do show is that a single significant study is
weak evidence of the size of an effect, and that the findings worth building a model on are those that
come back when the study is repeated.

\begin{measure}[An effect size as a correlation]
Different experiments measure different things: effort in a task, a price, a share of cooperators. To
compare them, a replication project converts each test statistic into a correlation coefficient $r$
between the treatment and the outcome. A difference of means $\Delta$ between two groups of equal size,
with a common standard deviation $s$, gives $r = \Delta/\sqrt{\Delta^2 + 4s^2}$. The conversion keeps
the sign and the strength of an effect but not its units: $r = 0.3$ in an effort task and $r = 0.3$ in
a market says nothing about which effect matters more for an economy.
\end{measure}

\section{Reference dependence and loss aversion}\label{sec:b2:behavioural-evidence:loss-aversion}

Expected utility values final positions: a \omterm{def:b2:expected-utility:lottery}{lottery} is good or bad according to the wealth it leaves.
Many choices look instead as if people valued changes, gains and losses measured from some point.

\begin{definition}[Reference point]\label{def:b2:behavioural-evidence:reference-point}
A \emph{reference point}\index{reference point} $z_0$ is the outcome against which a person codes
outcomes as gains ($z > z_0$) or losses ($z < z_0$): the status quo, an expectation, a target or the
outcome of others.
\end{definition}

\begin{definition}[Loss aversion]\label{def:b2:behavioural-evidence:loss-aversion}
A person whose preferences depend on a \omterm{def:b2:behavioural-evidence:reference-point}{reference point} shows \emph{loss aversion}\index{loss aversion}
if a loss reduces value by more than a gain of the same size raises it. With a \omterm{def:b2:behavioural-evidence:value-function}{value function} that is
linear on each side, $v(z - z_0) = z - z_0$ for gains and $\kappa(z - z_0)$ for losses, loss aversion
means $\kappa > 1$, and $\kappa$ is the \emph{loss-aversion coefficient}.
\end{definition}

The cleanest measurement uses \emph{mixed gambles}: a 50--50 chance of losing $z_-$ or gaining $z_+$. With
the linear \omterm{def:b2:behavioural-evidence:value-function}{value function}, the gamble is accepted when $\tfrac12 z_+ - \tfrac12\kappa z_- \ge 0$, that
is when $z_+/z_- \ge \kappa$. The smallest gain that makes a person accept, divided by the loss, measures
her $\kappa$. Expected utility predicts almost exactly one for small stakes: a utility of wealth that is
differentiable is nearly linear over a few hundred dalens, so any gamble with $z_+ > z_-$ should be
accepted. The calibration bound of \cref{prop:b2:expected-utility:calibration} says more: a person who
turns down small gambles with $z_+$ a little above $z_-$ at every level of wealth would, under expected
utility, have to turn down gambles that no one turns down.

\begin{record}[How large is loss aversion?]
Brown, Imai, Vieider and Camerer (2024) collected 607 estimates of the loss-aversion coefficient from
% ledger: b2:bivc2024
150 articles in economics, psychology, neuroscience and other fields. Their meta-analytic mean is
$1.955$, with a 95 per cent interval from $1.820$ to $2.102$: in the average study, a loss weighs about
twice as much as a gain of the same size. Estimates differ across designs and samples, and the meta-
analysis finds few study characteristics that predict them.
\end{record}

A coefficient of two is a summary across very different measurements: of mixed gambles, of choices
between insurance contracts, of prices in markets. It does not say what a person's \omterm{def:b2:behavioural-evidence:reference-point}{reference point} is,
and the estimate of $\kappa$ depends on it.

\begin{definition}[Endowment effect]\label{def:b2:behavioural-evidence:endowment-effect}
The \emph{endowment effect}\index{endowment effect} is a gap between the smallest price at which an
owner will sell a good (willingness to accept, WTA) and the largest price at which a non-owner will buy
it (willingness to pay, WTP), larger than income effects can explain.
\end{definition}

\omterm{def:b2:behavioural-evidence:loss-aversion}{Loss aversion} predicts the \omterm{def:b2:behavioural-evidence:endowment-effect}{endowment effect}. If giving up the good is a loss and acquiring it a gain,
and money is coded neutrally, a good worth $u$ in gains is sold at $\mathrm{WTA} = \kappa u$ and bought
at $\mathrm{WTP} = u$, a ratio of $\kappa$. If paying money is also felt as a loss, $\mathrm{WTP} =
u/\kappa$ and the ratio is $\kappa^2$. Under expected utility the two prices differ only by an income
effect, which is tiny for goods that cost a small part of wealth (\cref{exo:b2:behavioural-evidence:11}).

\begin{record}[Mugs]
Kahneman, Knetsch and Thaler (1990) gave coffee mugs at random to half the participants of an experiment
% ledger: b2:kkt1990
and then ran markets between owners and non-owners. Since the mugs were assigned at random, owners' and
non-owners' valuations come from the same distribution, so about half the mugs should change hands. In
every session, significantly fewer did: owners' willingness to accept greatly exceeded non-owners'
willingness to pay. In markets for tokens that the experimenters redeemed at a known value, the
predicted volume traded. Trading costs or misunderstanding the market would have reduced trade in tokens
too, so the authors concluded that they cannot explain the gap for mugs.
\end{record}

The token design is worth remembering: it holds the market institution and its costs fixed and changes
only whether the good has a use value of its own. It is the template for any test that separates a
preference from a cost of trading (\cref{exo:b2:behavioural-evidence:9}).

Two of the 18 replications of the opening chart tested reference dependence in the laboratory. Abeler,
% ledger: b2:camerer2016-abeler
Falk, Goette and Huffman (2011) found that participants worked longer at a dull task when their
expected earnings, the presumed \omterm{def:b2:behavioural-evidence:reference-point}{reference point}, were raised; the original effect size was $r = 0.18$
and the replication's $0.08$, not significant. Ericson and Fuster (2011) found that people asked a
% ledger: b2:camerer2016-ericson
higher price for a mug when they had been told they were more likely to receive it; $r = 0.21$ in the
original, $0.12$ and not significant in the replication. Both replications point in the original
direction with about half the effect. \omterm{def:b2:behavioural-evidence:loss-aversion}{Loss aversion} over the status quo is better established than the
claim that expectations set the \omterm{def:b2:behavioural-evidence:reference-point}{reference point}.

\section{Present bias in the field}\label{sec:b2:behavioural-evidence:present-bias}

\Cref{ch:b2:intertemporal-choice} built the quasi-hyperbolic model, in which a household weights every
future period by an extra factor $\varphi < 1$ relative to today, and found two pieces of field evidence.
% ledger: b2:dm2006
Gym members on flat monthly contracts paid more than \usd{17} per visit when a 10-visit pass cost
\usd{10} per visit, as if they expected to go more often than they did. And participants allocating dull
% ledger: b2:ans2013-beta
tasks across dates showed a present-bias factor of $0.900$ for effort, but $0.974$ for money, close to
no bias at all.

Two lessons carry over to the rest of this chapter. First, the evidence for \omterm{def:b2:intertemporal-choice:present-bias}{present bias} is stronger for
effort and consumption than for money, as the model predicts once money can be saved: shifting a payment
in time shifts nothing that is consumed. Second, the field evidence is indirect. A gym contract chosen
by someone who overestimates her future visits is consistent with \omterm{def:b2:intertemporal-choice:present-bias}{present bias} and naivety, but also
with poor information about one's own tastes. The model earns its place when the same parameter
predicts behaviour in another setting, here the demand for commitment, which appeared in the experiment
% ledger: b2:ans2013-commit
but vanished as soon as it had a price of 25 cents.

\section{Defaults}\label{sec:b2:behavioural-evidence:defaults}

Many decisions come with an option that takes effect if the person does nothing: a pension contribution
rate, a mobile-phone plan, an organ-donation status, the privacy settings of a website. Firms set some of
these defaults and governments others.

\begin{definition}[Default effect]\label{def:b2:behavioural-evidence:default-effect}
A \emph{default effect}\index{default effect} is a difference in the share of people who end up with an
option according to whether that option applies when they make no choice, the set of options and their
prices being the same.
\end{definition}

Under the standard model, a default can matter only through a cost of choosing. Suppose person $i$
gains $b_i$ from joining a pension plan, with $b_i$ normally distributed with mean $\mu$ and standard
deviation one across employees, and that acting on a form costs $\chi > 0$. Under opt-in, $i$ joins if
$b_i > \chi$; under opt-out, $i$ stays enrolled unless $b_i < -\chi$. Participation is
\[
P_{\text{in}} = \Phi(\mu - \chi), \qquad P_{\text{out}} = \Phi(\mu + \chi),
\]
where $\Phi$ is the standard normal distribution function. The gap rises with the cost $\chi$, and two
observed participation rates identify both $\mu$ and $\chi$ (\cref{pb:b2:behavioural-evidence:1}). The
cost must include more than the time to fill in a form: attention, procrastination of the kind modelled
in \cref{ch:b2:intertemporal-choice}, and the reading of the default as advice all raise the effective
$\chi$.

\begin{record}[Automatic enrolment]
Madrian and Shea (2000) studied a large American company that changed its retirement saving plan from
% ledger: b2:ms2000-participation
opt-in to automatic enrolment, leaving the plan's terms otherwise unchanged. At 3--15 months of tenure,
37 per cent of the employees hired just before the change participated, against 86 per cent of those
hired under automatic enrolment. The default contribution rate was 3 per cent of pay, all of it placed
% ledger: b2:ms2000-default
in a money-market fund. Of the employees hired under automatic enrolment, 61 per cent were at both
defaults, against 1 per cent of earlier cohorts; three quarters of the automatically enrolled
participants contributed 3 per cent, while three quarters of the others contributed 6 per cent or more.
\end{record}

\begin{omchart}{Retirement-plan participation and the default, at one American company, in per cent of
  employees with 3--15 months of tenure: hired just before automatic enrolment and hired under it.
  ``At the default'': contributing 3 per cent of pay, all in the money-market fund.}{Madrian and Shea
  (2000), NBER Working Paper 7682}
\begin{tikzpicture}
\begin{axis}[omrecord, ybar, bar width=22pt, xmin=0.4, xmax=2.6, ymin=0, ymax=100, xtick={1,2},
    xticklabels={Participating, At the default}, ylabel={Per cent of employees}, grid=none,
    ymajorgrids, width=0.62\linewidth, height=0.45\linewidth, nodes near coords,
    every node near coord/.append style={font=\footnotesize},
    legend style={at={(0.5,-0.2)}, anchor=north, font=\footnotesize},
    legend columns=2]
  \addplot[fill=omThm!60, draw=omThm] table[x=pos, y=window]
    {figdata/bachelor-2/05-behavioural-evidence/out/defaults.csv};
  \addplot[fill=omDef!70, draw=omDef] table[x=pos, y=new]
    {figdata/bachelor-2/05-behavioural-evidence/out/defaults.csv};
  \legend{Hired before, Hired under automatic enrolment}
\end{axis}
\end{tikzpicture}
\end{omchart}

Both bars move. A default changes who participates, and it also changes what participants do: the
default rate of 3 per cent was low against what earlier employees chose for themselves, so the same
default that raised participation pulled many savers' contributions down. Which effect dominates for
total saving is a question of numbers (\cref{pb:b2:behavioural-evidence:1}, Part III). The cost model
also says who is affected: those whose benefit $b_i$ lies between $-\chi$ and $\chi$, who follow
whichever default they are given. Under opt-in, some of them stay out although $b_i > 0$; under
opt-out, some stay in although $b_i < 0$. Neither default puts everyone where they would choose to be
if choosing were free, and the record cannot say which errors cost more. A firm that sets a default may
also have interests of its own, such as the cost of its matching contribution, and a government that
sets one may be wrong about what most people need; the evidence measures the effect of a default, not
which default is best.

\section{Limited attention}\label{sec:b2:behavioural-evidence:attention}

The standard model assumes that a consumer uses every price she faces. A tax added at the till, a fee
in small print, the price of ink for a printer bought cheaply: each is a price, and each may be paid
attention to less than the price on the shelf.

\begin{definition}[Salience]\label{def:b2:behavioural-evidence:salience}
A component of a price is \emph{salient}\index{salience} if it is displayed where the choice is made.
A consumer with \emph{limited attention} to a non-salient component $\tau$ chooses as if the price were
$p + \theta\tau$, with $0 \le \theta \le 1$: $\theta = 1$ is full attention, $\theta = 0$ none.
\end{definition}

With linear demand $q = a - \beta_1(p + \theta\tau)$, where $\beta_1$ is the slope a regression of
quantity on price would estimate, the effect of the tax on the quantity bought is
$-\beta_1\theta\tau$ instead of $-\beta_1\tau$. Posting the tax-inclusive price raises $\theta$ to one
and lowers demand; comparing the two reveals $\theta$ (\cref{exo:b2:behavioural-evidence:7}).

\begin{record}[Taxes not on the price tag]
Chetty, Looney and Kroft (2009) posted tax-inclusive price tags on some products in a grocery store, in
% ledger: b2:clk2009
a field experiment. Demand for those products fell by 8 per cent, although the tax paid at the till
was unchanged. In a second test, they found that increases in taxes included in posted prices reduced
alcohol consumption more than increases in taxes added at the register.
\end{record}

Limited attention has consequences that the standard model does not have, as Chetty, Looney and Kroft
% ledger: b2:clk2009
show. The share of a tax borne by buyers then depends on which side of the market is legally liable for it, and a tax that changes no one's
behaviour can still reduce welfare, since consumers who ignore it at the till pay it out of the
spending they had planned for other goods. Firms that sell to inattentive buyers have a reason to
shroud prices, putting part of the price where it is not seen. In models of shrouded prices, a
competitor who reveals the hidden part may gain little, if revealing it mainly teaches customers to buy
the cheap base product from the rival that still shrouds; then shrouding survives competition. Governments that require prices to be posted inclusive of taxes and fees change $\theta$ by
law.

\section{Findings that did not hold up}\label{sec:b2:behavioural-evidence:failures}

Seven of the 18 replications of the opening chart did not find a significant effect. They were findings
% ledger: b2:camerer2016-chen, b2:camerer2016-duffy, b2:camerer2016-huck
that a group identity made \omterm{def:b2:behavioural-evidence:salience}{salient} raises effort in a coordination game ($r = 0.84$ in the original,
$0.17$ in the replication); that monetary exchange raises efficiency over gift exchange in a small
% ledger: b2:camerer2016-ifcher, b2:camerer2016-kuziemko
economy ($0.76$, then $-0.12$); that deferred pay raises effort ($0.83$, then $0.37$); that a positive
mood makes people more patient ($0.28$, then $-0.01$); that people near the bottom of a ranking take
risks to avoid last place ($0.28$, then $-0.12$); and the two findings on \omterm{def:b2:behavioural-evidence:reference-point}{reference points} of the
previous section.

Three points keep such a list in proportion. A non-significant replication is evidence of a smaller
effect, not proof that the effect is zero: the replication of deferred pay found $r = 0.37$, a large
% ledger: b2:camerer2016-huck
effect by most standards, but not significant in a sample of that size. The originals that failed
% ledger: b2:camerer2016-predictors
were mostly those with the higher $p$-values and the smaller samples, as
\cref{prop:b2:behavioural-evidence:filter} leads one to expect: only one of the eight originals with a
$p$-value below $0.01$ failed. And the project's own point of comparison, a replication of 100
% ledger: b2:camerer2016-rpp
psychology studies, found a significant effect in the original direction for 36 per cent of those that
had reported one, so the issue is one of method, not of one field. For economics, the consequence is practical: a model
fitted to a single laboratory finding should be treated as provisional until the finding has been
reproduced.

\section{Prospect theory in the model frame}\label{sec:b2:behavioural-evidence:prospect-theory}

\omterm{def:b2:behavioural-evidence:prospect-theory}{Prospect theory} gathers reference dependence, \omterm{def:b2:behavioural-evidence:loss-aversion}{loss aversion} and a distortion of probabilities into one
model of choice under risk.

\begin{definition}[Value function (prospect theory)]\label{def:b2:behavioural-evidence:value-function}
A \emph{value function}\index{value function (prospect theory)} assigns a value $v(z - z_0)$ to each
outcome relative to the \omterm{def:b2:behavioural-evidence:reference-point}{reference point}, with $v(0) = 0$, $v$ increasing, concave for gains, convex for
losses and steeper for losses. This book uses the form $v(x) = x^{\eta}$ for $x \ge 0$ and $v(x) =
-\kappa(-x)^{\eta}$ for $x < 0$, with $0 < \eta \le 1$ and $\kappa > 1$.
\end{definition}

\begin{definition}[Probability weighting function]\label{def:b2:behavioural-evidence:weighting}
A \emph{probability weighting function}\index{probability weighting function} $\omega : [0, 1] \to [0,
1]$ is increasing with $\omega(0) = 0$ and $\omega(1) = 1$; it replaces each probability by a decision
weight. This book uses $\omega(p) = p^{d}/(p^{d} + (1 - p)^{d})^{1/d}$ with $0 < d \le 1$, which is the
identity when $d = 1$.
\end{definition}

\begin{definition}[Prospect theory]\label{def:b2:behavioural-evidence:prospect-theory}
Under \emph{prospect theory}\index{prospect theory}, a person facing a \omterm{def:b2:expected-utility:lottery}{lottery} with at most one gain
$z_1 > z_0$ with probability $p_1$ and one loss $z_2 < z_0$ with probability $p_2$ values it at
\[
V = \omega(p_1)\,v(z_1 - z_0) + \omega(p_2)\,v(z_2 - z_0),
\]
and chooses the \omterm{def:b2:expected-utility:lottery}{lottery} with the highest $V$. Versions for many outcomes apply the weighting function to
cumulative probabilities, separately for gains and for losses.
\end{definition}

The figures show the two functions with the illustrative values $\eta = 0.9$, $\kappa = 2$ and $d =
0.65$ used in this chapter's exercises. They are not estimates; the loss-aversion coefficient is close to
the meta-analytic mean of the record above.

\begin{omfigure}
\begin{tikzpicture}[econ]
\begin{axis}[omrecord, grid=none, axis lines=middle, xmin=-110, xmax=110, ymin=-140, ymax=80,
    xtick={-100,-50,50,100}, ytick={-100,-50,50}, xlabel={$z - z_0$}, ylabel={$v$},
    x label style={anchor=north west}, y label style={anchor=south east},
    tick label style={font=\scriptsize}, width=0.62\linewidth, height=0.52\linewidth]
  \addplot[omDef, very thick, domain=0:105, samples=60] {x^0.9};
  \addplot[omThm, very thick, domain=-105:0, samples=60] {-2*(-x)^0.9};
  \addplot[guide, domain=-105:105, samples=2] {x};
  \node[anchor=north west, font=\footnotesize, omDef] at (axis cs:55,30) {gains: $z^{0.9}$};
  \node[anchor=north west, font=\footnotesize, omThm] at (axis cs:-80,-118) {losses: $-2(-z)^{0.9}$};
  \node[anchor=south east, font=\footnotesize, black!60] at (axis cs:-45,-38) {$v = z - z_0$};
\end{axis}
\end{tikzpicture}

{\small A prospect-theory \omterm{def:b2:behavioural-evidence:value-function}{value function} with $\eta = 0.9$ and $\kappa = 2$ (illustrative values). It is
kinked at the \omterm{def:b2:behavioural-evidence:reference-point}{reference point}: the slope just below it is twice the slope just above. A loss of 100
lowers value by about 126, a gain of 100 raises it by about 63.}
\end{omfigure}

\begin{omfigure}
\begin{tikzpicture}[econ]
\begin{axis}[omrecord, grid=none, xmin=0, xmax=1, ymin=0, ymax=1, xtick={0,0.25,0.5,0.75,1},
    ytick={0,0.25,0.5,0.75,1}, xlabel={Probability $p$}, ylabel={Decision weight $\omega(p)$},
    width=0.55\linewidth, height=0.55\linewidth]
  \addplot[guide, domain=0:1, samples=2] {x};
  \addplot[omDef, very thick, domain=0.0005:0.9995, samples=120] {x^0.65/(x^0.65 + (1-x)^0.65)^(1/0.65)};
  \node[eqpoint] at (axis cs:0.359,0.359) {};
  \node[anchor=north west, font=\footnotesize] at (axis cs:0.37,0.33) {$\omega(p) = p$ at $p \approx 0.36$};
  \node[anchor=south west, font=\footnotesize, omDef] at (axis cs:0.02,0.42) {overweighted};
  \node[anchor=north west, font=\footnotesize, omDef] at (axis cs:0.6,0.45) {underweighted};
\end{axis}
\end{tikzpicture}

{\small A \omterm{def:b2:behavioural-evidence:weighting}{probability weighting function} with $d = 0.65$ (illustrative). Small probabilities get more
weight than their value, large ones less; the curve crosses the diagonal at about $0.36$. The weight of
a 1 per cent chance is $0.047$, almost five times the probability.}
\end{omfigure}

% code: exercises/solutions.py crossing, ex_fourfold
The weighting function explains the Allais pattern of \cref{ch:b2:expected-utility} without any
curvature of utility.

\begin{proposition}[Common ratio under prospect theory]\label{prop:b2:behavioural-evidence:common-ratio}
With the \omterm{def:b2:behavioural-evidence:reference-point}{reference point} at zero, $\eta = 0.9$ and $d = 0.65$, \omterm{def:b2:behavioural-evidence:prospect-theory}{prospect theory} ranks the lotteries of
\cref{prop:b2:expected-utility:common-ratio} as $\ell^{B'} \succ \ell^{A'}$ and $\ell^{C'} \succ
\ell^{D'}$, the pattern most people choose and expected utility excludes.
\end{proposition}

\begin{proof}
% code: exercises/solutions.py allais
$V(\ell^{A'}) = \omega(0.8)\,4000^{0.9} = 0.640 \times 1745 = 1117$ and $V(\ell^{B'}) = 3000^{0.9} = 1347$,
so $\ell^{B'}$ is preferred. $V(\ell^{C'}) = \omega(0.2)\,4000^{0.9} = 0.260 \times 1745 = 454$ and
$V(\ell^{D'}) = \omega(0.25)\,3000^{0.9} = 0.293 \times 1347 = 395$, so $\ell^{C'}$ is preferred. The
ratio $\omega(0.2)/\omega(0.25) = 0.886$ is larger than $\omega(0.8)/\omega(1) = 0.640$: scaling both
probabilities down by four makes the difference between them matter less, which is what the
\omterm{def:b2:expected-utility:independence-axiom}{independence axiom} forbids.
\end{proof}

Combined with a \omterm{def:b2:behavioural-evidence:value-function}{value function} that is concave for gains and convex for losses, the weighting function
also predicts a \emph{fourfold pattern} of attitudes to risk: \omterm{def:b2:expected-utility:risk-aversion}{risk aversion} for likely gains and for
unlikely losses (people insure), risk seeking for unlikely gains (people buy \omterm{def:b2:expected-utility:lottery}{lottery} tickets) and for
likely losses (people gamble to avoid a sure loss). \Cref{exo:b2:behavioural-evidence:10} derives it.

\begin{model}{Prospect theory}
\assumes Outcomes are valued as gains and losses from a \omterm{def:b2:behavioural-evidence:reference-point}{reference point}; losses weigh more than gains
($\kappa > 1$); sensitivity diminishes away from the \omterm{def:b2:behavioural-evidence:reference-point}{reference point} ($\eta < 1$); probabilities are
replaced by decision weights that overweight small probabilities.
\predicts The Allais pattern; the \omterm{def:b2:behavioural-evidence:endowment-effect}{endowment effect}, with a WTA--WTP ratio of $\kappa$ or $\kappa^2$; the
rejection of small mixed gambles with positive expected value; the fourfold pattern of risk attitudes;
choices that change when the same outcomes are described as gains or as losses.
\tested Laboratory gambles, in hundreds of studies (the meta-analysis above); the mugs experiment;
choices of insurance \omterm{def:b2:insurance-and-risk-sharing:deductible}{deductibles}; field effects of \omterm{def:b2:behavioural-evidence:reference-point}{reference points}, such as the default in a pension
plan.
\breaks The theory does not say where the \omterm{def:b2:behavioural-evidence:reference-point}{reference point} is: the status quo, an expectation or a
target can each fit a given choice, and a model that chooses its \omterm{def:b2:behavioural-evidence:reference-point}{reference point} after seeing the data
can fit almost anything. Laboratory findings on expectations as \omterm{def:b2:behavioural-evidence:reference-point}{reference points} replicate weakly
(\cref{sec:b2:behavioural-evidence:loss-aversion}), and experienced traders show little \omterm{def:b2:behavioural-evidence:endowment-effect}{endowment
effect} (\cref{sec:b2:behavioural-evidence:experience}).
\end{model}

\section{Experienced traders and high stakes}\label{sec:b2:behavioural-evidence:experience}

Most of the evidence above comes from students choosing over small sums for an hour. Two objections
follow, and both can be tested. Markets give experience and feedback, and people may learn to avoid
costly anomalies; and stakes in real decisions are larger.

\begin{record}[Traders and the endowment effect]
List (2003) studied traders in two real marketplaces, combining field observation with an
% ledger: b2:list2003
experiment. As market experience increased, individual behaviour converged to the prediction of the
standard model, and the experimental test found that market experience played a significant role in
eliminating the \omterm{def:b2:behavioural-evidence:endowment-effect}{endowment effect}.
\end{record}

Experience may teach a trader that goods bought are goods to be sold again, so that parting with them is
no loss. The model frame then asks: is a \omterm{def:b2:behavioural-evidence:reference-point}{reference point} something a person learns? If so, \omterm{def:b2:behavioural-evidence:loss-aversion}{loss
aversion} is real but its domain is narrower than the laboratory suggests, and its effects should be
largest for goods held for use, by people with little practice at trading them.

\begin{debate}[Do taxi drivers stop at a daily income target?]
A driver who chooses her hours freely each day faces a wage that varies from day to day. The standard
model of labour supply predicts that she works more on high-wage days, since an hour of leisure costs
more then. A driver with a daily income target and \omterm{def:b2:behavioural-evidence:loss-aversion}{loss aversion} below it does the opposite: she
reaches the target sooner on a good day and goes home. Camerer, Babcock, Loewenstein and Thaler (1997)
% ledger: b2:cblt1997
estimated the daily wage elasticity of hours for New York cab drivers in three samples and found it
significantly negative in two, at about $-1$ for inexperienced drivers. Farber (2015) used the records of
% ledger: b2:farber2015
all New York taxi trips of 2009--2013, more than 700 million trips by about 62,000 drivers, and found
that drivers worked more when earnings opportunities rose, whether the rise was anticipated or not; he
found reference dependence not to be important for daily hours. The questions still open are how to measure the
wage, since a wage computed as a day's earnings divided by its hours falls mechanically whenever hours
are overstated, and whether a minority of drivers, or of days, does stop at a target.
\end{debate}

\Cref{exo:b2:behavioural-evidence:12} works the target model. Its prediction is precise: hours fall
with the wage, with an elasticity of $-1$, but only over a band of wages within which the target is
reached at the margin; outside the band, hours rise with the wage, as in the standard model. A test of
the model is therefore a test of where most driver-days lie, and that is what the larger data set
settled for New York.

The broad lesson of the chapter fits in the model frame. \omterm{def:b2:behavioural-evidence:loss-aversion}{Loss aversion} over the status quo, the effect
of defaults and the neglect of \omterm{def:b2:behavioural-evidence:salience}{non-salient} prices have been reproduced in many settings, in the field and
with real stakes. \omterm{def:b2:intertemporal-choice:present-bias}{Present bias} is clearer for effort than for money. Several laboratory findings have
not replicated, and experience removes at least one well-known \omterm{def:b2:asset-prices:anomaly}{anomaly}. Models built on these findings
are tools like the others in this book: each predicts something the standard model does not, and each
is kept as long as the prediction survives.

\section{Exercises}\label{sec:b2:behavioural-evidence:exercises}

\begin{exercise}[$\star$]\label{exo:b2:behavioural-evidence:1}
Read the opening chart. How many replications found a smaller effect than the original, and how many an
effect of the opposite sign? Compute the \omterm{def:b2:behavioural-evidence:replication-rate}{replication rate}. Why can a \omterm{def:b2:behavioural-evidence:replication-rate}{replication rate} of 61 per cent
coexist with every original effect being real?
\end{exercise}

\begin{exercise}[$\star$]\label{exo:b2:behavioural-evidence:2}
A person with a linear \omterm{def:b2:behavioural-evidence:value-function}{value function} and $\kappa = 2$ is offered a 50--50 gamble: lose \dal{100} or win
\dal{150}. Find its value. Does she accept? What is the smallest gain at which she accepts?
\end{exercise}

\begin{exercise}[$\star$]\label{exo:b2:behavioural-evidence:3}
A mug is worth $u = 5$ dalens of gain to anyone who gets it, and $\kappa = 2$. Find WTA and WTP if only
the mug is subject to \omterm{def:b2:behavioural-evidence:loss-aversion}{loss aversion}, then if paying money is also felt as a loss. What ratio WTA/WTP does
each version predict?
\end{exercise}

\begin{exercise}[$\star$]\label{exo:b2:behavioural-evidence:4}
Describe an observation that would show a person's rejection of a small mixed gamble is not due to \omterm{def:b2:behavioural-evidence:loss-aversion}{loss
aversion}, and one that would show a gap between WTA and WTP is not an \omterm{def:b2:behavioural-evidence:endowment-effect}{endowment effect}.
\end{exercise}

\begin{exercise}[$\star\star$]\label{exo:b2:behavioural-evidence:5}
With $v(x) = x^{0.9}$ for gains, $v(x) = -2(-x)^{0.9}$ for losses and no probability weighting, value the
50--50 gamble that wins \dal{1000} or loses \dal{400}. Is it accepted? What is its expected value, and
what would expected utility with a nearly linear utility predict?
\end{exercise}

\begin{exercise}[$\star\star$]\label{exo:b2:behavioural-evidence:6}
A firm of 2,000 employees moves its pension plan from opt-in to automatic enrolment. Using the
participation and default rates of the automatic-enrolment record, predict the number of participants
before and after, and the number at the default. If the average salary is \dal{30000} and the default
rate 3 per cent, what do the employees at the default contribute in a year? What does the record say
about the contributions of those who would have joined anyway?
\end{exercise}

\begin{exercise}[$\star\star$]\label{exo:b2:behavioural-evidence:7}
Demand for a good is $q = 1600 - 40(p + \theta\tau)$, with $p = 20$ and a sales tax $\tau = 2$ added at
the till. When the tax-inclusive price is posted, $\theta = 1$. Find demand when the tax is posted.
Posting reduces demand by 8 per cent, as in the grocery experiment. Find demand without posting and the
attention weight $\theta$ it implies.
\end{exercise}

\begin{exercise}[$\star\star$]\label{exo:b2:behavioural-evidence:8}
In Table~S1 of the replication project, the reference-point studies of Abeler et al.\ and of Ericson and
Fuster had original effect sizes $0.18$ and $0.21$ and replication effect sizes $0.08$ and $0.12$; the
relative effect sizes reported are $0.43$ and $0.58$. Check the ratios. Neither replication was
significant. Write two sentences that report these results accurately, and one that would overstate them.
\end{exercise}

\begin{exercise}[$\star\star$]\label{exo:b2:behavioural-evidence:9}
Design an experiment that separates \omterm{def:b2:behavioural-evidence:loss-aversion}{loss aversion} from transaction costs as the cause of low trading
volume, in the spirit of the mugs and tokens design. State what each hypothesis predicts in each
treatment.
\end{exercise}

\begin{exercise}[$\star\star\star$]\label{exo:b2:behavioural-evidence:10}
Use $\omega(p)$ with $d = 0.65$ and a \omterm{def:b2:behavioural-evidence:value-function}{value function} linear on each side. (a) Compute $\omega(0.01)$,
$\omega(0.05)$, $\omega(0.5)$, $\omega(0.9)$ and $\omega(0.99)$. (b) Compare a 1 per cent chance of
\dal{1000} with \dal{10} for sure, and a 99 per cent chance of \dal{1000} with \dal{990} for sure. (c) Do
the same for losses of the same sizes. (d) State the fourfold pattern your answers display and the
everyday choices it is used to explain.
\end{exercise}

\begin{exercise}[$\star\star\star$]\label{exo:b2:behavioural-evidence:11}
A person has utility $\ln y + \gamma_m\,\mathbb{1}[\text{mug}]$ over money $y$ and owning a mug, with
$\gamma_m > 0$, and money $m$. (a) Show that $\mathrm{WTP} = m(1 - e^{-\gamma_m})$ and $\mathrm{WTA} =
m(e^{\gamma_m} - 1)$, so that WTA/WTP $= e^{\gamma_m}$. (b) For a mug for which she would pay \dal{5}
out of $m = \dal{30000}$, find the ratio. (c) Prove more generally that if utility is continuously differentiable in money and owning the good
does not change the marginal utility of money, WTA/WTP tends to one as the value of the good tends to
zero, so that expected utility cannot produce a ratio near two for small stakes.
\end{exercise}

\begin{exercise}[$\star\star\star$]\label{exo:b2:behavioural-evidence:12}
A driver earns $w$ dalens per hour and chooses her hours $h$ each day. Her daily income is $m = wh$ and
her utility is $m - \bar m - \nu h^2/2$ if $m \ge \bar m$ and $\kappa(m - \bar m) - \nu h^2/2$ if $m <
\bar m$, with target $\bar m = 200$, $\nu = 4$ and $\kappa = 2$. (a) Show that she stops exactly at the
target when $w$ lies between $20$ and $\sqrt{800} = 28.28$, and find her hours at $w = 16$, $20$, $25$,
$28$ and $32$. (b) Find the wage elasticity of hours inside and outside the band. (c) What does the
standard model ($\kappa = 1$) predict? (d) What pattern in data on many driver-days would reject the
target model?
\end{exercise}

\section{Problem: Brenmoor's pension plan}\label{sec:b2:behavioural-evidence:problem}

\begin{problem}[{Weekend problem --- Brenmoor's town council, in Brenland, an invented economy, measures
its residents' loss aversion, chooses between opt-in and opt-out for its employees' pension plan, and
predicts their choice of insurance}]
\label{pb:b2:behavioural-evidence:1}
The council surveys 300 residents. Each is offered, for real money, 50--50 gambles that lose \dal{100}
or win $z_+$, for $z_+ = 100, 125, \dots, 400$. The table gives the percentage who accept (invented
data, \texttt{gamble\_survey.csv}). Round $\kappa$ to three decimals, normal quantiles and the
parameters of the default model to four, and percentages to one decimal.

\begin{center}
\small
\begin{tabular}{l*{7}{r}}
\toprule
Gain $z_+$ (dalens) & 100 & 125 & 150 & 175 & 200 & 225 & 250 \\
% figdata: gamble_survey (accept by G)
Accept (per cent) & 1.3 & 8.7 & 20.3 & 35.3 & 50.0 & 63.7 & 73.7 \\
\midrule
Gain $z_+$ (dalens) & 275 & 300 & 325 & 350 & 375 & 400 & \\
Accept (per cent) & 82.0 & 87.7 & 90.7 & 93.7 & 96.0 & 97.3 & \\
\bottomrule
\end{tabular}
\end{center}

\medskip
\noindent\textbf{Part I --- Loss aversion from a survey.}
\begin{enumerate}
  \item What share of the residents should accept the gamble with $z_+ = 125$ under expected utility
        with a smooth utility of wealth? Why?
  \item What share rejects it in the survey?
  \item Find the smallest gain at which at least half accept, and the loss-aversion coefficient of the
        median resident with a linear \omterm{def:b2:behavioural-evidence:value-function}{value function}.
  \item With $v(x) = x^{0.9}$ for gains and $-\kappa(-x)^{0.9}$ for losses, find the median resident's
        $\kappa$.
  \item Compare these estimates with the meta-analytic interval of the record.
  \item A resident turns down the gamble with $z_+ = 125$ at every level of wealth. Under expected
        utility, using the bound of \cref{exo:b2:expected-utility:10}, $(z_- + z_+)z_+/(z_+ - z_-)$, which
        50--50 gambles must she also turn down, whatever their gain?
\end{enumerate}

\medskip
\noindent\textbf{Part II --- An endowment test.}
\begin{enumerate}[resume]
  \item The council gives a town mug to half of a class of residents at random. Owners ask \dal{8} on
        average; non-owners offer \dal{3.5}. Find the ratio.
  \item Find the $\kappa$ implied if only the mug is subject to \omterm{def:b2:behavioural-evidence:loss-aversion}{loss aversion}, and if money is too.
  \item Which version agrees with Part I?
  \item What would a token market in the same class tell the council?
\end{enumerate}

\medskip
\noindent\textbf{Part III --- Opt-in or opt-out?} Each employee gains $b_i$ from joining the plan, normal
with mean $\mu$ and standard deviation one; acting costs $\chi$. Take the participation rates of the
automatic-enrolment record (37 and 86 per cent) to calibrate $\chi$.
\begin{enumerate}[resume]
  \item Write the two participation rates as functions of $\mu$ and $\chi$.
  \item Solve for $\chi$ and $\mu$ from the record's rates.
  \item Brenmoor's current opt-in plan has 45 per cent participation. With the same $\chi$, find
        Brenmoor's $\mu$.
  \item Predict participation under opt-out, and the gap between the two defaults.
  \item Under opt-in, participants choose 6 per cent of pay. Under opt-out, 61 per cent of employees stay
        at the default of 3 per cent, as in the record, and the other participants choose 6 per cent.
        Find the average contribution per employee under each default.
  \item Under opt-out, what share of employees participates although $b_i < 0$?
  \item Under opt-in, what share stays out although $b_i > 0$?
  \item Does the model say which default is better? What would it need to know?
\end{enumerate}

\medskip
\noindent\textbf{Part IV --- A prediction for insurance.} Residents insure their homes. A \omterm{def:b2:insurance-and-risk-sharing:deductible}{deductible} of
\dal{500} costs \dal{100} a year more than a \omterm{def:b2:insurance-and-risk-sharing:deductible}{deductible} of \dal{1000}; the chance of a claim in a year is
5 per cent. Assume that residents code the premium as a cost, not a loss, but an uninsured part of a
claim as a loss, weighted by $\kappa$ and by $\omega(p)$ with $d = 0.65$, and value linear on each side.
\begin{enumerate}[resume]
  \item How much does the low \omterm{def:b2:insurance-and-risk-sharing:deductible}{deductible} save in expectation each year? What does expected utility, with
        a smooth utility and a wealth of \dal{50000}, predict?
  \item Find $\omega(0.05)$ and the $\kappa$ above which a resident buys the low \omterm{def:b2:insurance-and-risk-sharing:deductible}{deductible}.
  \item If $\kappa$ across residents is lognormal with median 2 and standard deviation of $\ln\kappa$
        equal to $0.35$ (the survey's design), what share buys the low \omterm{def:b2:insurance-and-risk-sharing:deductible}{deductible}?
  \item What data from the town's insurer would test this prediction, and what result would reject it?
  \item What else, apart from \omterm{def:b2:behavioural-evidence:loss-aversion}{loss aversion} and probability weighting, could make residents buy low
        \omterm{def:b2:insurance-and-risk-sharing:deductible}{deductibles}?
  \item State the result: the estimated loss-aversion coefficient, the predicted participation gap
        between the two defaults, and the predicted share buying the low \omterm{def:b2:insurance-and-risk-sharing:deductible}{deductible}.
\end{enumerate}
\end{problem}
